\paragraph{Theorem F.2.26 (Proved: the complete rational remaining $(5,6)$ locus).}
The irreducible rational rank-five pairs in the second row
of (F.S1), with square-zero rank-two cusp, are classified
up to simultaneous rational conjugacy by
\[
 \{[r:s:t]\in\mathbb P^2(\mathbb Q):
             D(r,s,t)\ne0,\ Q_-(r,s,t)\ne0\}
       \ \sqcup\ \mathbb A^1(\mathbb Q),
 \qquad Q_-=r^2-rs+st,
 \tag{F.R10}
\]
where $D$ is the same cubic as in (F.R2).
The exceptional line again lies over $[4:1:3]$.
Thus the two rational spectral loci left by F.2.25 both
have complete explicit normal forms, including their
entire rational exceptional loci.

\begin{proof}
Use precisely the space $\mathbb Q\oplus\mathbb Q(\zeta_5)$,
matrices $A,R,X$ and projective plane normalization in
(F.R3)--(F.R5).  The centralizer and image-plane arguments
in F.2.20 require only irreducibility and the two-dimensional
$-1$ eigenspace of $C$; both hold here.
Consequently every pair has a unique projective image-plane
parameter, with only scalar stabilizer.
The unknown $Y$ now satisfies
\[
 YX=0,\qquad YRX=I_2,\qquad
 \operatorname{tr}(YAX)=-2,\qquad
 \operatorname{tr}(YA^{-1}X)=2.
 \tag{F.R11}
\]
Its coefficient matrix is unchanged from (F.R6), so has
determinant $D/4$.
The two traces and the $-1$ eigenspace imply
\[
 \chi_C=(x-1)(x+1)^2(x^2+x+1).
\]
Its other three roots are distinct, and the $-1$ eigenspace
has its full multiplicity, so $C^6=I$.
The cusp $XY$ still has square zero and rank two.
For $D\ne0$, the ten linear equations give the unique
explicit solution.

The two first-column adjugate identities become
\[
 D y_{10}=2Q_-Q_2,
 \qquad D y_{20}=2(r-s-t)Q_-,
 \tag{F.R12}
\]
with the same $Q_2$ as in (F.R7).
For $D\ne0$, these entries vanish simultaneously exactly
when $Q_-=0$.
Indeed, if $Q_-\ne0$ they would instead force
$r=s+t$ and $Q_2=0$, and (F.R8) would give $D=0$.
The invariant-covector argument of F.2.20 is unchanged.
Thus the pair is rationally irreducible exactly off the
conic $Q_-=0$ in the generic locus.

No nonzero rational triple satisfies $D=Q_-=0$.
If $s=0$, the conic forces $r=0$, and then $D=7t^3$.
If $s\ne0$, rescale to $s=1$ and put
$r=u$, $t=u-u^2$.  Direct substitution gives
\[
 D(u,1,u-u^2)=
 -(7u^2-11u+7)(u^4-u^3+u^2-u+1).
 \tag{F.R13}
\]
The first factor has discriminant $-75$ and positive
leading coefficient.  The second is strictly positive,
as the identity
\[
 u^4-u^3+u^2-u+1
 =\left(u^2-\frac u2\right)^2
  +\frac34\left(u-\frac23\right)^2+\frac23
\]
shows.  In particular neither has a rational zero.
On the exceptional cubic the necessary identities (F.R12)
therefore force the same two points $[3:2:1]$ and $[4:1:3]$
as before.  The first is again impossible because
$[X\mid RX]$ has rank less than four.
At the second, the complete solution is
\[
 Y_z=\begin{pmatrix}
 4z+26/5&0&-2z-2&-2z-1/5&2z+1\\
 2z-2/5&0&-z&2/5-z&z
 \end{pmatrix},\qquad z\in\mathbb Q.
 \tag{F.R14}
\]
Its first-column entries cannot vanish simultaneously,
so each pair on this line is irreducible.
There are no further identifications, by the scalar
stabilizer argument already proved in F.2.20.
This establishes the claimed bijection.

The exact program is
\texttt{fourfolds\_relative\_six\_probe.py}.
It checks the modified characteristic polynomials, finite-order
identities, first-column adjugate expressions, conic
substitution and entire exceptional line symbolically.
Its auxiliary trace-parameter computations for the other
repeated-negative spectrum are recorded as well, although
that spectrum is now excluded integrally by F.2.25.
No finite sample is used to infer a rational classification.
The finite lattice and extension algorithm of F.2.22 applies
to every point of (F.R10); it does not by itself decide
which rational parameters survive the remaining integral
homology tests.
\end{proof}
